\chapter{Worked Example: A Lattice Boltzmann Kernel}\label{r01:ch:worked-example}

This chapter combines the pieces of Part~I into one realistic kernel: a fused
stream-and-collide step of the D2Q9 lattice Boltzmann method (LBM) with the BGK
collision operator. This is the kind of kernel lbmpy~\cite{lbmpy} generates and waLBerla
runs. It uses multi-component fields, a memory layout chosen for vectorization,
subexpressions, common subexpression elimination, a scalar parameter, and two targets.
The complete script is \file{examples/ex13_lbm.py}.

\section{The model}

In the LBM, every cell stores $Q=9$ \term{populations} $f_q$, one per lattice velocity
$\mathbf c_q$. One time step consists of
\begin{align*}
  \text{streaming:}\quad & \tilde f_q(\mathbf x) = f_q(\mathbf x - \mathbf c_q),\\
  \text{collision:}\quad & f_q^{\text{new}}(\mathbf x) = \tilde f_q - \omega\,\bigl(\tilde f_q - f_q^{\text{eq}}(\rho, \mathbf u)\bigr),\\
  \text{with}\quad & \rho = \sum_q \tilde f_q,\qquad \rho\,\mathbf u = \sum_q \mathbf c_q \tilde f_q,\\
  & f_q^{\text{eq}} = w_q\,\rho\,\bigl(1 + 3\,\mathbf c_q\!\cdot\!\mathbf u + \tfrac92 (\mathbf c_q\!\cdot\!\mathbf u)^2 - \tfrac32\,\mathbf u\!\cdot\!\mathbf u\bigr).
\end{align*}

\begin{figure}[htbp]
  \centering
  % D2Q9: lattice velocities (numbered as in examples/ex13_lbm.py) and pull streaming.
\begin{tikzpicture}[font=\footnotesize]
  \def\vel{0/0/0,1/1/0,2/0/1,3/-1/0,4/0/-1,5/1/1,6/-1/1,7/-1/-1,8/1/-1}
  % ---- (a) velocity set -----------------------------------------------------------
  \begin{scope}[x=1.15cm,y=1.15cm]
    \foreach \i in {-1,0,1} {
      \draw[gridline] (\i,-1.35) -- (\i,1.35); \draw[gridline] (-1.35,\i) -- (1.35,\i);
      \foreach \j in {-1,0,1} \fill[psgray] (\i,\j) circle (1.2pt);
    }
    \foreach \q/\cx/\cy in \vel {
      \ifnum\q=0
        \fill[psblue!70] (0,0) circle (2.6pt);
        \node[celltext] at (0.33,-0.15) {0};
      \else
        \draw[-{Stealth[length=2.2mm]},thick,psblue!80!black] (0,0) -- (\cx*0.9,\cy*0.9);
        \node[celltext,fill=white,inner sep=1pt] at (\cx*1.17,\cy*1.17) {\q};
      \fi
    }
    \node[font=\small] at (0,1.95) {(a) velocities $\mathbf c_q$, $q = 0\dots8$};
    \node[align=center,celltext] at (0,-2.0)
      {$w_0 = \frac49$,\ \ $w_{1\dots4} = \frac19$,\ \ $w_{5\dots8} = \frac1{36}$};
  \end{scope}

  % ---- (b) pull streaming -----------------------------------------------------------
  \begin{scope}[shift={(6.3cm,0)},x=1.15cm,y=1.15cm]
    \fill[psblue!25] (-0.5,-0.5) rectangle (0.5,0.5);
    \foreach \i in {-1.5,-0.5,0.5,1.5} {
      \draw[gridline] (\i,-1.5) -- (\i,1.5); \draw[gridline] (-1.5,\i) -- (1.5,\i);
    }
    \foreach \q/\cx/\cy in \vel {
      \ifnum\q=0 \else
        \draw[-{Stealth[length=2.2mm]},thick,psorange!90!black]
          (-\cx*0.95,-\cy*0.95) -- (-\cx*0.3,-\cy*0.3);
        \node[celltext,fill=white,inner sep=0.5pt] at (-\cx*1.05,-\cy*1.05) {$f_{\q}$};
      \fi
    }
    \node[celltext] at (0,0) {$\mathbf x$};
    \node[font=\small] at (0,1.95) {(b) pull: $\tilde f_q(\mathbf x) = f_q(\mathbf x - \mathbf c_q)$};
    \node[align=center,celltext] at (0,-2.0)
      {e.g.\ $q=1$ ($\mathbf c_1 = (1,0)$): \texttt{src[-1,0](1)}};
  \end{scope}
\end{tikzpicture}

  \caption[The D2Q9 model and pull streaming]{The D2Q9 model. (a) The nine lattice velocities with their weights, numbered as in the example script. (b) Pull streaming: cell $\mathbf x$ fetches population $q$ from the neighbour at $\mathbf x-\mathbf c_q$, which in \ps{} is the field access \texttt{src[-cx, -cy](q)}.}
  \label{r01:fig:d2q9}
\end{figure}

The ``pull'' formulation (\cref{r01:fig:d2q9}) reads the neighbours' populations and writes
only the own cell.
This makes the kernel free of loop-carried dependencies, as \ps{} requires.

\section{The symbolic description}

Both population arrays are 2D fields with an index dimension of size 9. The layout
\code{fzyx} (structure of arrays) makes $x$ the contiguous coordinate. That is the layout
SIMD needs, because neighbouring cells of the same population are adjacent in memory.
\pysnippet{ex13_lbm.model}
The rules use only SymPy. Density and velocity are subexpressions, and every main
assignment writes one population. \code{ps.simp.sympy_cse} then removes the many
repeated subterms:
\pysnippet{ex13_lbm.rules}
\outsnippet{ex13_lbm.ac-size}

\section{Generating and running}

The same assignment collection produces a portable kernel and a vectorized kernel for
the best SIMD target of the host:
\pysnippet{ex13_lbm.kernels}
\outsnippet{ex13_lbm.stats}
The arrays are created with the matching memory order (population slowest, $x$
fastest), including one ghost layer:
\pysnippet{ex13_lbm.run}
\outsnippet{ex13_lbm.check}
Both kernels match an independent NumPy implementation of the same step
(\code{np.roll} for streaming, \code{einsum} for the moments) on all interior cells.

The beginning of the generated portable kernel shows how the CSE subexpressions
(\code{xi_*}) and the density became \code{const} locals. With the \code{fzyx} layout the
innermost loop runs over \code{ctr_0} ($x$), and the population index enters as
\code{q * _stride_src_2}:
\csnippet[nolol=false,caption={The generated D2Q9 kernel (first 40 lines).},lastline=40,basicstyle=\ttfamily\scriptsize]{ex13_lbm.scalar}

\section{From here to production}

\begin{itemize}
  \item \textbf{Real models.} lbmpy derives such update rules (moment- and
        cumulant-based collision operators, forcing schemes, boundary conditions) and
        optimizes them symbolically before handing them to \ps.
  \item \textbf{Real applications.} In waLBerla the same \code{Kernel} object is
        embedded into a C++ sweep through sweepgen and pystencils-sfg
        (\cref{r01:sec:external-use}). It then runs on distributed, block-structured grids
        with MPI.
  \item \textbf{This workspace.} \file{python/learn-lb-pystencil/lb-pystencils.py} is a
        complete D2Q9 lid-driven-cavity and Taylor--Green simulation built this way.
        It has command-line switches for the target (\code{generic}/\code{rvv}),
        precision, compiler and fused or split kernels, and is validated against the
        NumPy reference in \file{python/learn-lb-numpy}.
        \file{scripts/lbm_rvv_check.py} compares the RVV and generic collide kernels.
        The harness numbers quoted in \cref{r01:ch:rvv-backend} were measured with
        \file{lb-pystencils.py}, the waLBerla numbers with the
        \file{scripts/walberla_*.sh} benchmarks.
\end{itemize}

\section{Further reading}
\begin{itemize}
  \item The upstream user manual and developer reference~\cite{pystencils-docs}, whose
        sources are in \file{submodules/pystencils/docs/source}. The backend walkthrough
        (\file{backend/backend_walkthrough.md}) and the GPU code generation notes
        (\file{backend/gpu_codegen.md}) complement Part~II of this manual.
  \item The tests in \file{submodules/pystencils/tests}: \file{nbackend/} for the
        backend, \file{kernelcreation/} for end-to-end behaviour, \file{jit/} for the
        JIT. They are the most precise specification of the current behaviour.
  \item The workspace notes in \file{docs/notes} and \file{docs/notes2} on the RVV
        backend, its performance, and the waLBerla code generation flow.
\end{itemize}
